% CATWALK: A Stream Cipher Based on Arnold's Cat Map
% in a Coupled Lattice Sponge Construction
%
% First complete draft for IACR ePrint submission.
% Status: research cipher, self-reviewed, awaiting independent cryptanalysis.

\documentclass[runningheads]{llncs}

\usepackage[utf8]{inputenc}
\usepackage[T1]{fontenc}
\usepackage{amsmath,amssymb,amsfonts}
\usepackage{algorithm}
\usepackage{algpseudocode}
\usepackage{tikz}
\usepackage{pgfplots}
\pgfplotsset{compat=1.18}
\usepackage{booktabs}
\usepackage{multirow}
\usepackage{hyperref}
\usepackage{xcolor}
\usepackage{enumitem}
\usepackage{todonotes}

% Security claim labels
\newcommand{\proof}[1]{\textsc{[Proof]} #1}
\newcommand{\conjecture}[1]{\textsc{[Conjecture]} #1}
\newcommand{\empirical}[1]{\textsc{[Empirical]} #1}

% Common notation
\newcommand{\Ztwo}{\mathbb{Z}/2^{64}\mathbb{Z}}
\newcommand{\catmap}{\mathbf{A}}
\newcommand{\coupmat}{\mathbf{C}}
\newcommand{\state}{\mathbf{s}}
\newcommand{\GOLDEN}{\texttt{0x9E3779B97F4A7C15}}

\begin{document}

\title{CATWALK: A Stream Cipher Based on Arnold's Cat Map\\in a Coupled Lattice Sponge Construction}

\author{Anonymous}
\institute{Independent Research}

\maketitle

%% ============================================================
%% ABSTRACT
%% ============================================================
\begin{abstract}
Can spatiotemporal chaos serve as a cryptographic primitive?
We present CATWALK, an authenticated encryption scheme whose core permutation
is a Coupled Map Lattice (CML) wrapped in a Bertoni~et~al.\ sponge framework.
The local map is Arnold's Cat Map---a linear toral automorphism with proven
Anosov hyperbolicity---operating on 8 adjacent pairs of a 16-site lattice
over~$\Ztwo$.  Unlike prior chaos-based ciphers, CATWALK uses
exact integer arithmetic (no floating point), a standard sponge construction
with formal capacity analysis, and a local map whose hyperbolic properties are
provable, not assumed.  A principled eigenvalue analysis of the coupling
topology discovered and fixed a structural singularity, leading to a 5-term
coupling architecture with distances $\{1,3,7,11\}$ that is fully invertible
over~$\Ztwo$ (trivial kernel, $\det = -33075$, $\min|\lambda_k| = 1.259$).
PractRand validation to 1~TB per seed (397 tests per seed at the 1~TB checkpoint;
full multi-seed validation in progress), raw
permutation validation at~32~GB without the output finalizer, and
reduced-round analysis establishing a $4\times$ security margin all support
the empirical quality of the construction.  CATWALK is a research cipher: it
has not received independent cryptanalysis, and the core security conjecture---that the CML-Sponge permutation is indistinguishable from a random
permutation---remains unproven.  We present the design, analysis, and open
problems to invite community scrutiny.
\end{abstract}

\keywords{stream cipher \and authenticated encryption \and coupled map lattice \and Arnold's Cat Map \and sponge construction \and chaos-based cryptography}

%% ============================================================
%% 1. INTRODUCTION
%% ============================================================
\section{Introduction}
\label{sec:intro}

Chaos theory offers a compelling intuition for cryptographic design:
a chaotic dynamical system exhibits sensitive dependence on initial conditions
(small input changes produce large, unpredictable output changes), topological
mixing (any open set eventually maps onto the entire phase space), and dense
periodic orbits (the system visits every neighborhood infinitely often).
These properties mirror the classical Shannon requirements of confusion and
diffusion~\cite{arnold1968ergodic}.

\paragraph{Why previous chaos ciphers failed.}
Despite this intuition, chaos-based cryptography has a troubled history.
Baptista's seminal proposal~\cite{baptista1998cryptography} and many subsequent
designs~\cite{kocarev2001chaos,fridrich1998symmetric,chen2004symmetric}
suffer from three fundamental problems identified by Li~et~al.~\cite{li2001problems}:
\begin{enumerate}[nosep]
\item \textbf{Floating-point non-determinism.}
  Continuous chaotic maps require real-number arithmetic.
  Finite-precision IEEE~754 floating-point introduces rounding errors that
  accumulate unpredictably across platforms, destroying cross-platform
  determinism and reducing effective dynamical complexity.
\item \textbf{No formal framework.}
  Most chaos ciphers are \emph{ad hoc} stream generators without a
  principled construction (no sponge, no Feistel, no SPN).
  This makes formal security analysis impossible: there is no capacity
  bound, no indifferentiability proof, and no well-defined security game.
\item \textbf{Weak or absent mixing analysis.}
  The chaotic properties of the continuous map do not automatically transfer
  to the discretized version.  Complement symmetry, short periods on finite
  fields, and low-order bit bias have broken multiple proposals.
\end{enumerate}

\paragraph{What CATWALK does differently.}
CATWALK addresses each of these problems with three specific design decisions:
\begin{enumerate}[nosep]
\item \textbf{Exact integer arithmetic.}
  Arnold's Cat Map is the matrix $\begin{psmallmatrix}1&1\\1&2\end{psmallmatrix}$
  acting on~$(\Ztwo)^2$ via wrapping additions---no floating point, no
  approximation, no precision loss.  The map is natively integer.
\item \textbf{Sponge framework.}
  The CML permutation is wrapped in a standard Bertoni--Daemen--Peeters--Van~Assche
  sponge~\cite{bertoni2007sponge,bertoni2008indifferentiability},
  providing rate/capacity separation, domain-separated absorption, and
  a formal security bound of $q^2/2^c$ (where $c = 512$ bits is the capacity).
\item \textbf{Provable hyperbolicity.}
  Arnold's Cat Map is an Anosov diffeomorphism on the real torus $\mathbb{T}^2$,
  with Lyapunov exponent $\ln((3+\sqrt{5})/2) \approx 0.9624$.
  This is a \emph{proven} mathematical property, not a numerical observation.
  The map has no complement symmetry (Section~\ref{sec:construction}),
  its unique fixed point at $(0,0)$ is rendered negligible by Weyl counter
  injection, and its determinant~1 guarantees bijectivity over any ring.
\end{enumerate}

\paragraph{Design-by-analysis methodology.}
The design of CATWALK's coupling topology illustrates a principled
cryptographic design process.
A geometric analysis of the lattice dynamics revealed that the original
coupling distances $\{1,7,8\}$ produced a \emph{singular} coupling matrix---the
polynomial $p(x) = 1+x+x^7+x^8$ has $p(-1) = 0$, meaning the alternating
Fourier mode is in the kernel, destroying the sponge security guarantee
(Section~\ref{sec:coupling}).
This discovery motivated a systematic eigenvalue analysis of all feasible
coupling topologies:  an exhaustive search over all $\binom{15}{4} = 1365$
distance combinations identified 160~valid 5-term candidates, from which
$\{1,3,7,11\}$ was selected for its optimal properties
($\det(\coupmat) = -33075$, odd; $\min|\lambda_k| = 1.259$; trivial kernel;
all-prime distances).
The singularity discovery and its resolution through principled search is
the central methodological contribution of this paper.

\paragraph{Paper contributions.}
\begin{enumerate}[nosep]
\item A complete specification of the CML-Sponge AEAD construction---precise
  enough to reimplement from scratch (Section~\ref{sec:construction}).
\item A comprehensive coupling matrix analysis, including the discovery and
  resolution of a structural singularity through exhaustive polynomial root
  analysis (Section~\ref{sec:coupling}).
\item A security analysis with honest labeling of every claim as \textsc{[Proof]},
  \textsc{[Conjecture]}, or \textsc{[Empirical]} (Section~\ref{sec:security}).
\item Statistical validation at 1~TB with PractRand, raw permutation validation
  at 32~GB without the output finalizer, and reduced-round analysis
  establishing a $4\times$ conservative security margin (Section~\ref{sec:validation}).
\item An explicit enumeration of open problems inviting independent
  cryptanalysis (Section~\ref{sec:openproblems}).
\end{enumerate}

\paragraph{Paper roadmap.}
Section~\ref{sec:prelim} covers mathematical preliminaries.
Section~\ref{sec:construction} gives the full construction specification.
Section~\ref{sec:coupling} presents the coupling matrix analysis.
Section~\ref{sec:security} develops the security argument.
Section~\ref{sec:validation} reports statistical validation results.
Section~\ref{sec:performance} discusses performance.
Section~\ref{sec:related} surveys related work.
Section~\ref{sec:openproblems} enumerates open problems.
Section~\ref{sec:conclusion} concludes.

%% ============================================================
%% 2. PRELIMINARIES
%% ============================================================
\section{Preliminaries}
\label{sec:prelim}

\subsection{Coupled Map Lattices}

A \emph{Coupled Map Lattice} (CML), introduced by Kaneko~\cite{kaneko1984cml,kaneko1985cml},
is a dynamical system defined on a regular lattice of~$N$ sites,
where each site $i$ carries a state $s_i$ in some phase space~$\mathcal{S}$.
The system evolves in discrete time through two alternating steps:
\begin{enumerate}[nosep]
\item \textbf{Local map:} each site applies a map $f: \mathcal{S} \to \mathcal{S}$
  independently.
\item \textbf{Coupling:} each site's state is updated by a weighted sum of its
  own state and those of its neighbors.
\end{enumerate}
For a ring lattice of $N$~sites with coupling function $g$ and distances
$\mathcal{D} = \{d_1, \ldots, d_k\}$:
\[
  s_i^{(t+1)} = g\!\left(\{f(s_{(i+d) \bmod N}^{(t)})\}_{d \in \mathcal{D} \cup \{0\}}\right).
\]
CMLs exhibit \emph{spatiotemporal chaos}: the interaction of local chaotic dynamics
with spatial coupling produces complex, high-dimensional behavior that is
resistant to low-dimensional approximation.

\subsection{Sponge Constructions}

The sponge construction, introduced by Bertoni~et~al.~\cite{bertoni2007sponge},
is a mode of operation for a fixed-width permutation~$\pi: \{0,1\}^b \to \{0,1\}^b$.
The $b$-bit state is partitioned into a \emph{rate} $r$ and a \emph{capacity} $c$
($b = r + c$).  During \emph{absorption}, input blocks are XOR'd into the rate
portion and $\pi$ is applied.  During \emph{squeezing}, rate bits are read as
output and $\pi$ is applied.

The fundamental security result~\cite{bertoni2008indifferentiability} states that
a sponge with an ideal permutation is indifferentiable from a random oracle,
with distinguishing advantage bounded by:
\[
  \mathrm{Adv} \leq \frac{q^2}{2^c}
\]
where $q$ is the number of queries.  This bound requires the capacity to be
\emph{effectively} $c$~bits---i.e., the permutation must not have structural
weaknesses that reduce the capacity below~$c$.

Keccak~\cite{keccak2011}, the SHA-3 standard, is the canonical sponge-based hash.
SpongeWrap extends the sponge to authenticated encryption by absorbing ciphertext
(not plaintext) for authentication, providing native Encrypt-then-MAC semantics.

\subsection{Arnold's Cat Map}

Arnold's Cat Map~\cite{arnold1968ergodic} is the linear toral automorphism
defined by the matrix:
\[
  \catmap = \begin{pmatrix} 1 & 1 \\ 1 & 2 \end{pmatrix}, \qquad
  \begin{pmatrix} x' \\ y' \end{pmatrix}
  = \catmap \begin{pmatrix} x \\ y \end{pmatrix} \pmod{2^{64}}.
\]
Computed in symplectic order (update $x$ first, then use updated $x$ to compute $y$): $x' = x + y$, then $y' = x' + y = x + 2y$.

\paragraph{Key properties.}
\begin{itemize}[nosep]
\item \textbf{Area-preserving:} $\det(\catmap) = 2 - 1 = 1$.
  The map is a bijection on $(\Ztwo)^2$ and on $(\mathbb{Z}/m\mathbb{Z})^2$ for any~$m$.
\item \textbf{Hyperbolic (Anosov):} eigenvalues $\lambda_{\pm} = (3 \pm \sqrt{5})/2
  \approx \{2.618, 0.382\}$, both real, product~1, neither a root of unity.
  The map has stable and unstable manifolds filling~$\mathbb{R}^2$.
\item \textbf{Lyapunov exponent:} $\Lambda = \ln\lambda_+ = \ln((3+\sqrt{5})/2) \approx 0.9624$.
  For any $2 \times 2$ integer matrix with $\det = 1$ and trace~3,
  the characteristic equation $\lambda^2 - 3\lambda + 1 = 0$ gives
  eigenvalues $(3 \pm \sqrt{5})/2$ uniquely, so the Lyapunov exponent
  is the same for all such matrices.
\item \textbf{No complement symmetry:}
  $\catmap(M-x, M-y) \neq \catmap(x,y)$ and $\neq (M - \catmap(x,y))$ for
  generic inputs, where $M = 2^{64} - 1$.  See Section~\ref{sec:construction}.
\item \textbf{Unique fixed point:} $(0,0)$ is the only fixed point.
\item \textbf{Symplectic:} $\catmap \in \mathrm{Sp}(2, \mathbb{Z})$.
  The map preserves the symplectic form $dx \wedge dy$.
\end{itemize}

\subsection{Notation}

\begin{table}[h]
\centering
\caption{Notation used throughout this paper.}
\label{tab:notation}
\begin{tabular}{@{}ll@{}}
\toprule
Symbol & Meaning \\
\midrule
$N = 16$ & Number of lattice sites \\
$r = 512$ & Rate (bits); sites $s[0..7]$ \\
$c = 512$ & Capacity (bits); sites $s[8..15]$ \\
$R = 8$ & Rounds per permutation \\
$\catmap$ & Arnold's Cat Map matrix $\begin{psmallmatrix}1&1\\1&2\end{psmallmatrix}$ \\
$\coupmat$ & $16 \times 16$ circulant coupling matrix \\
$p(x)$ & Coupling polynomial $1 + x + x^3 + x^7 + x^{11}$ \\
$\omega$ & Primitive 16th root of unity $e^{2\pi i/16}$ \\
$\lambda_k$ & $k$-th eigenvalue of $\coupmat$: $\lambda_k = p(\omega^k)$ \\
$\oplus$ & Bitwise XOR \\
$\boxplus$ & Wrapping addition mod~$2^{64}$ \\
$\boxtimes$ & Wrapping multiplication mod~$2^{64}$ \\
\bottomrule
\end{tabular}
\end{table}

%% ============================================================
%% 3. THE CATWALK CONSTRUCTION
%% ============================================================
\section{The CATWALK Construction}
\label{sec:construction}

This section specifies the CML-Sponge AEAD construction precisely enough
to reimplement from scratch.

\subsection{State Layout}

The permutation operates on a 1024-bit state consisting of:
\begin{itemize}[nosep]
\item \textbf{Lattice:} 16 words $s[0], \ldots, s[15]$, each $\in \Ztwo$.
\item \textbf{Weyl counter:} one word $\mathrm{ctr} \in \Ztwo$.
\end{itemize}
The lattice is partitioned as:
\begin{itemize}[nosep]
\item \textbf{Rate} ($r = 512$ bits): sites $s[0..7]$.
  XOR'd with input during absorb; read during squeeze.
\item \textbf{Capacity} ($c = 512$ bits): sites $s[8..15]$.
  Never directly output; carries secret state across permutations.
\end{itemize}
Initial state: all sites zero, counter zero.

\subsection{One Permutation Round}

Each round applies four steps in sequence.  See Algorithm~\ref{alg:round}.

\begin{algorithm}[t]
\caption{One CML-Sponge round.  All arithmetic is wrapping mod $2^{64}$.}
\label{alg:round}
\begin{algorithmic}[1]
\Require State $s[0..15]$, counter $\mathrm{ctr}$
\Statex \textbf{Step 1 --- Weyl counter injection:}
\State $\mathrm{ctr} \gets \mathrm{ctr} \boxplus \mathrm{GOLDEN}$
  \Comment{$\mathrm{GOLDEN} = \lfloor \varphi \cdot 2^{64} \rfloor = \GOLDEN$}
\For{$i = 0$ \textbf{to} $15$}
  \State $s[i] \gets s[i] \boxplus \mathrm{rotl}(\mathrm{ctr},\, \mathrm{ROT}[i])$
    \Comment{ROT = first 16 primes $\geq 3$}
\EndFor
\Statex \textbf{Step 2 --- Arnold's Cat Map (snapshot):}
\For{$k = 0$ \textbf{to} $7$}
  \State $m[2k] \gets s[2k] \boxplus s[2k+1]$
  \State $m[2k+1] \gets m[2k] \boxplus s[2k+1]$
    \Comment{Symplectic order: $y' = x' + y$}
\EndFor
\Statex \textbf{Step 3 --- CML 5-term additive coupling:}
\For{$i = 0$ \textbf{to} $15$}
  \State $s[i] \gets m[i] \boxplus m[(i{+}1)\bmod 16] \boxplus m[(i{+}3)\bmod 16]
    \boxplus m[(i{+}7)\bmod 16] \boxplus m[(i{+}11)\bmod 16]$
\EndFor
\Statex \textbf{Step 4 --- Multiplicative mixing:}
\For{$k = 0$ \textbf{to} $7$}
  \State $s[2k{+}1] \gets s[2k{+}1] \boxtimes (s[2k] \mathbin{|} 1)$
    \Comment{$|\,1$ ensures odd multiplier}
\EndFor
\end{algorithmic}
\end{algorithm}

\paragraph{Step 1: Weyl counter injection.}
The Weyl sequence with $\mathrm{GOLDEN} = \lfloor \varphi \cdot 2^{64} \rfloor$
($\varphi = (1+\sqrt{5})/2$) is a nothing-up-my-sleeve constant with optimal
equidistribution.  The 16 rotation amounts $\mathrm{ROT} = [3, 5, 7, 11, 13, 17,
19, 23, 29, 31, 37, 41, 43, 47, 53, 61]$ are the first 16 primes $\geq 3$,
all coprime to~64.  Each site receives a distinct rotation of the counter,
breaking translational symmetry across pairs before the Cat Map is applied.

\paragraph{Step 2: Arnold's Cat Map.}
The matrix $\catmap = \begin{psmallmatrix}1&1\\1&2\end{psmallmatrix}$
is applied to the 8~adjacent pairs $(s[0],s[1]), (s[2],s[3]), \ldots, (s[14],s[15])$.
All 8~pairs are written into a snapshot array~$m[]$ before the coupling step,
preserving single-step semantics.  This pairing matches Step~4's multiplicative
mixing pairs, creating a coherent two-layer transformation per pair.

\paragraph{Step 3: CML additive coupling.}
The coupling distances $\{1, 3, 7, 11\}$ produce the polynomial
$p(x) = 1 + x + x^3 + x^7 + x^{11}$ with $p(1) = 5$ (odd).
Each site receives the sum of itself and four neighbors.
Full analysis is in Section~\ref{sec:coupling}.

\paragraph{Step 4: Multiplicative mixing.}
$s[2k+1] \gets s[2k+1] \boxtimes (s[2k] \mathbin{|} 1)$.
The OR with~1 ensures the multiplier is always odd, hence invertible
mod~$2^{64}$.
\proof{Step~4 is a bijection on each pair $(s[2k], s[2k+1])$ for every value of~$s[2k]$,
since odd integers are units in~$\Ztwo$.}
This is the sole nonlinear step in the round function---Steps 1--3 are affine
over~$\Ztwo$.

\subsection{Permutation and Output Finalizer}

The permutation $\pi$ applies $R = 8$ rounds of Algorithm~\ref{alg:round}.
Before output, the Stafford Mix13 bijective finalizer~\cite{stafford2011}
is applied to each rate word:
\begin{align*}
  x &\gets x \oplus (x \gg 30), \quad
  x \gets x \boxtimes \texttt{0xBF58476D1CE4E5B9}, \\
  x &\gets x \oplus (x \gg 27), \quad
  x \gets x \boxtimes \texttt{0x94D049BB133111EB}, \\
  x &\gets x \oplus (x \gg 31).
\end{align*}
\proof{Mix13 is a bijection on~$\Ztwo$ (composition of xor-shifts with odd multiplications).
It satisfies the strict avalanche criterion: every output bit depends on every input bit.}

\empirical{Raw permutation output (pre-Mix13) passes PractRand at 32~GB
on 3 seeds with zero failures (Section~\ref{sec:validation}), confirming
that Mix13 is defense-in-depth, not load-bearing.}

\subsection{Sponge Operations}

\paragraph{Absorb.}
Keccak-style multi-rate padding: $M \| D \| \texttt{0x00}^* \| \texttt{0x80}$,
where $D$ is the domain byte and padding brings the total to a multiple of
64~bytes.  Each 64-byte block is XOR'd into the rate (sites~0--7, little-endian)
and the permutation $\pi$ is applied.

\proof{The padding is injective: distinct $(M, D)$ pairs produce distinct padded blocks.}

\paragraph{Squeeze.}
Apply Mix13 to each rate word, read 64~bytes (little-endian), apply~$\pi$.

\paragraph{Domain separation.}
Five domain bytes ensure that absorb calls for different roles cannot produce
identical state transitions (Table~\ref{tab:domains}).

\begin{table}[h]
\centering
\caption{Domain separation bytes.}
\label{tab:domains}
\begin{tabular}{@{}lll@{}}
\toprule
Domain & Value & Usage \\
\midrule
\texttt{DOMAIN\_KEY} & \texttt{0x01} & Key absorption \\
\texttt{DOMAIN\_IV}  & \texttt{0x02} & IV/nonce absorption \\
\texttt{DOMAIN\_AAD} & \texttt{0x03} & Associated data (authenticated, not encrypted) \\
\texttt{DOMAIN\_CT}  & \texttt{0x04} & Ciphertext absorption during AEAD \\
\texttt{DOMAIN\_TAG} & \texttt{0x05} & Tag finalization (empty absorb) \\
\bottomrule
\end{tabular}
\end{table}

\subsection{AEAD Construction}

The AEAD construction follows SpongeWrap semantics.
Algorithm~\ref{alg:aead-enc} specifies encryption.
Decryption is identical except: keystream is generated
first, then ciphertext is absorbed (not plaintext),
then XOR'd to recover plaintext.  The tag is verified
before any plaintext is returned.

\begin{algorithm}[t]
\caption{CATWALK AEAD encryption.}
\label{alg:aead-enc}
\begin{algorithmic}[1]
\Require Key $K$ (32 bytes), nonce $N$ (16 bytes),
         header $H$, plaintext $P$
\Ensure Ciphertext $C$, tag $T$
\State $\sigma \gets \mathsf{CipherInit}(K, N)$
  \Comment{Absorb $K$ with domain 0x01, $N$ with domain 0x02}
\State $\mathsf{Absorb}(\sigma, H, \texttt{0x03})$
  \Comment{Header authenticated, not encrypted}
\ForAll{64-byte blocks $P_i$ of $P$}
  \State $K_i \gets \mathsf{Squeeze}(\sigma, 64)$
    \Comment{Squeeze keystream (with Mix13)}
  \State $C_i \gets P_i \oplus K_i$
  \State $\mathsf{Absorb}(\sigma, C_i, \texttt{0x04})$
    \Comment{Absorb ciphertext, not plaintext}
\EndFor
\State $\mathsf{Absorb}(\sigma, \varepsilon, \texttt{0x05})$
  \Comment{Domain-separated finalisation}
\State $T \gets \mathsf{Squeeze}(\sigma, 32)$
  \Comment{32-byte tag}
\State \Return $(C, T)$
\end{algorithmic}
\end{algorithm}

\paragraph{Domain separation.}
Five domain bytes separate the five absorption phases:
\texttt{0x01}~(key), \texttt{0x02}~(IV/nonce),
\texttt{0x03}~(AAD/header), \texttt{0x04}~(ciphertext),
\texttt{0x05}~(tag finalisation).
The Keccak-style multi-rate padding
$M \| D \| \texttt{0x00}^* \| \texttt{0x80}$
makes the padded encoding injective: distinct
$(M, D)$ pairs produce distinct padded blocks,
preventing cross-phase collisions.

\proof{Both encrypt and decrypt absorb the \emph{ciphertext} with
domain~\texttt{0x04}, so sponge state evolution is identical.
The tag is bound to: key, nonce, all AAD, and all ciphertext---providing
native Encrypt-then-MAC semantics.}

\subsection{Key Derivation}

\begin{enumerate}[nosep]
\item Argon2id~\cite{rfc9106}: password $+$ (16-byte random salt $\|$ 8-byte timestamp)
  $\to$ 32-byte master key.  Default: $m = 256$~MB, $t = 4$ iterations, $p = 1$ lane.
\item BLAKE3~\cite{blake3} domain derivation:
  $\mathrm{cipher\_key} = \mathrm{BLAKE3\_derive\_key}(\text{``catwalk.v9.cipher''},\, \mathrm{master\_key})$.
\item Sponge init: absorb cipher\_key, then absorb nonce.
\end{enumerate}

Minimum accepted parameters on decryption: $m \geq 64$~MB, $t \geq 2$.
This prevents crafted headers from bypassing the KDF cost.

\subsection{Parameters Summary}

\begin{table}[h]
\centering
\caption{CATWALK parameters.}
\label{tab:params}
\begin{tabular}{@{}lll@{}}
\toprule
Parameter & Value & Rationale \\
\midrule
Lattice sites & 16 & Power of 2; rich coupling topology \\
Site width & 64 bits & Matches CPU word width \\
Total state & 1024 + 64 bits & 16 sites + Weyl counter \\
Rate / Capacity & 512 / 512 bits & Full capacity (no correction) \\
Rounds & 8 & $4\times$ diffusion margin \\
Coupling distances & $\{1,3,7,11\}$ & 5-term; $|\det|=33075$ (odd); $\min|\lambda|=1.259$ \\
GOLDEN & $\GOLDEN$ & $\lfloor \varphi \cdot 2^{64} \rfloor$ \\
ROT & First 16 primes $\geq 3$ & All coprime to 64 \\
Output finalizer & Stafford Mix13 & Bijective whitening layer \\
Tag size & 32 bytes & 128-bit forgery resistance (birthday) \\
Nonce size & 16 bytes & 128-bit; collision at $\sim\!2^{64}$ encryptions \\
KDF & Argon2id & Memory-hard; RFC~9106 \\
\bottomrule
\end{tabular}
\end{table}

%% ============================================================
%% 4. COUPLING MATRIX ANALYSIS
%% ============================================================
\section{Coupling Matrix Analysis}
\label{sec:coupling}

This section presents the systematic analysis that discovered a structural
singularity in the original coupling topology and resolved it through
exhaustive polynomial root analysis.  This design-by-analysis process is the
central methodological contribution of this paper.

\subsection{Circulant Structure}

The coupling step (Step~3) applies a $16 \times 16$ circulant matrix $\coupmat$
over~$\Ztwo$ whose first row encodes the coupling polynomial
$p(x) = \sum_{d \in \mathcal{D} \cup \{0\}} x^d$.
The eigenvalues of a circulant over~$\mathbb{C}$ are
$\lambda_k = p(\omega^k)$ for $k = 0, \ldots, 15$,
where $\omega = e^{2\pi i/16}$ is the primitive 16th root of unity.

\subsection{The Original \texorpdfstring{$\{1,7,8\}$}{1,7,8} Design and Its Failure}

The initial CATWALK design used distances $\{1,7,8\}$, giving:
\[
  p(x) = 1 + x + x^7 + x^8.
\]
Evaluating at $x = -1$ (the 16th root of unity with $k = 8$):
\[
  p(-1) = 1 + (-1) + (-1)^7 + (-1)^8 = 1 - 1 - 1 + 1 = 0.
\]
\proof{$\lambda_8 = p(-1) = 0$, so $\coupmat$ is singular over~$\mathbb{C}$ (and over every ring).}

The null vector $\mathbf{v} = (1, -1, 1, -1, \ldots, 1, -1) \in \ker(\coupmat)$
means that two states differing by~$\mathbf{v}$ produce \emph{identical} coupling
output.  This collapses one dimension of the state space per round, destroying
the sponge security guarantee: the effective capacity is reduced by at least
one bit per coupling application.

\paragraph{Why this matters.}
The sponge bound $\mathrm{Adv} \leq q^2/2^c$ assumes the permutation operates
on a full $b$-bit state.  If the coupling step has a non-trivial kernel, the
permutation's effective state dimension is reduced, and the capacity bound no
longer holds at its stated value.  For the $\{1,7,8\}$ design, the rank
deficiency is catastrophic: it provides an algebraic distinguisher for the
permutation, breaking the random-permutation assumption entirely.

\subsection{Fix Methodology: Polynomial Root Analysis}

The singularity at $p(-1) = 0$ arises because the distances $\{1,7,8\}$
include an even distance (8), so $(-1)^8 = +1$, which can cancel with the
odd-distance terms.  This observation leads to a clear design requirement:

\begin{quote}
\emph{Choose coupling distances such that $p(\omega^k) \neq 0$ for all
$k = 0, \ldots, 15$.  Equivalently, the coupling polynomial must have no roots
among the 16th roots of unity.}
\end{quote}

Additionally, for the coupling to be invertible over~$\Ztwo$ (preserving full
sponge capacity), we require $\det(\coupmat)$ to be \emph{odd}---since
$\gcd(\det(\coupmat), 2^{64}) = 1$ if and only if $\det(\coupmat)$ is odd.

\subsection{3-Term Coupling: Evaluated and Rejected}

A 3-term polynomial $p(x) = 1 + x^{d_1} + x^{d_2}$ has $p(1) = 3$ (odd),
which immediately fixes the invertibility issue.  However, all 100 valid 3-term
candidates have $\geq 3$ unit-magnitude eigenvalues ($|\lambda_k| = 1$ for at
least 3 values of~$k$).  A unit-magnitude eigenvalue means the coupling step
performs \emph{zero mixing work} on that Fourier mode---it is an isometry
providing no diffusion.

\proof{For any 3-term unit-coefficient polynomial with 2~non-zero exponents chosen
from $\{1, \ldots, 15\}$, at least 3 of the 16~eigenvalue magnitudes equal~1.
This is a structural impossibility to avoid given the degree constraints of
3-term circulants on a 16-site ring.}

\subsection{4-Term \texorpdfstring{$\{1,5,11\}$}{1,5,11}: Partial Fix}

The intermediate design used distances $\{1,5,11\}$:
$p(x) = 1 + x + x^5 + x^{11}$, $p(1) = 4$ (even).
\[
  \det(\coupmat) = -1088 = -2^6 \times 17 \quad (\text{even}).
\]
Since $\det(\coupmat)$ is even, $\gcd(1088, 2^{64}) = 64 \neq 1$, and $\coupmat$
is \emph{not} invertible over~$\Ztwo$.  The kernel has 4~elements:
$\{0,\, 2^{62}\mathbf{1},\, 2^{63}\mathbf{1},\, 3 \cdot 2^{62}\mathbf{1}\}$,
causing a 2-bit effective capacity reduction (510~bits instead of~512).

Additionally, $\min|\lambda_k| = 0.765$ (at $k = 2, 6, 10, 14$), meaning
25\% of Fourier modes are \emph{contracted} by the coupling---they receive
less energy after coupling than before.

\subsection{5-Term Exhaustive Search}

For the 5-term polynomial $p(x) = 1 + x^{d_1} + x^{d_2} + x^{d_3} + x^{d_4}$
with $1 \leq d_1 < d_2 < d_3 < d_4 \leq 15$, there are $\binom{15}{4} = 1365$
candidate distance sets.  We evaluated all 1365 against four requirements:

\begin{enumerate}[nosep,label=R\arabic*.]
\item Non-singular over $\mathbb{C}$: $|p(\omega^k)| > 0$ for all~$k$.
\item Odd determinant: $\det(\coupmat)$ odd (invertible over~$\Ztwo$).
\item No unit-magnitude eigenvalues: $||p(\omega^k)| - 1| > 10^{-6}$ for $k \neq 0$.
\item Full 16-site diffusion in $\leq 4$ rounds.
\end{enumerate}

\textbf{160 candidates} satisfy all four requirements.  Of these, 134 achieve
2-round full diffusion, and 40 achieve the maximum eigenvalue margin
$\min|\lambda_k| = 1.259$.  Among the top-tier candidates, 6 have all-odd
distances.

\subsection{Selection of \texorpdfstring{$\{1,3,7,11\}$}{1,3,7,11}}

The recommended distance set $\{1,3,7,11\}$ was selected from the 6 all-odd
top-tier candidates based on:
\begin{itemize}[nosep]
\item \textbf{All-odd distances} $\Rightarrow$ $p(-1) = 1-1-1-1-1 = -3 \neq 0$
  (algebraic guarantee).
\item \textbf{All-prime distances} (1, 3, 7, 11): nothing-up-my-sleeve character.
\item \textbf{Well-distributed} across $\{1, \ldots, 15\}$ (not clustered).
\item \textbf{Shares} $d_1 = 1$ and $d_4 = 11$ with the prior design, making the
  change targeted.
\end{itemize}

\begin{table}[h]
\centering
\caption{Direct comparison of coupling designs.}
\label{tab:coupling-comparison}
\begin{tabular}{@{}lccc@{}}
\toprule
Property & 3-term (best) & 4-term $\{1,5,11\}$ & 5-term $\{1,3,7,11\}$ \\
\midrule
$p(1)$ & 3 (odd) & 4 (even) & 5 (odd) \\
Invertible over $\Ztwo$? & Yes & No & \textbf{Yes} \\
Unit-magnitude modes & $\geq 3$ & 0 & \textbf{0} \\
$\min|\lambda_k|$ & 1.000 & 0.765 & \textbf{1.259} \\
$|\det(\coupmat)|$ & varies & 1088 & \textbf{33075} \\
Kernel & trivial & 4 elements & \textbf{trivial} \\
Effective capacity & 512 bits & 510 bits & \textbf{512 bits} \\
Diffusion (rounds) & varies & 2 & \textbf{2} \\
\bottomrule
\end{tabular}
\end{table}

\subsection{Proofs of Non-Singularity}

\proof{(Over $\mathbb{C}$.)  All 16 eigenvalues $\lambda_k = p(\omega^k)$ are
non-zero.  The minimum magnitude is $\min_k |\lambda_k| = 1.259$ (at
$k = 1, 7, 9, 15$).  Direct computation confirms $|p(\omega^k)| > 1$ for all~$k$.
The determinant is $\det(\coupmat) = \prod_{k=0}^{15} \lambda_k = -33075$.}

\proof{(Over $\mathbb{Z}$.)  The same eigenvalue computation applies: the
characteristic polynomial of~$\coupmat$ over~$\mathbb{Q}$ has no zero among
the 16th roots of unity, so $\coupmat$ is non-singular over~$\mathbb{Q}$ and
hence over~$\mathbb{Z}$.}

\proof{(Over $\Ztwo$.)  $\det(\coupmat) = -33075 = -3^3 \times 5^2 \times 7^2$
(odd).  $\gcd(33075, 2^{64}) = 1$, so $\coupmat$ is fully invertible
over~$\Ztwo$.  The kernel is trivial: $\ker(\coupmat) = \{\mathbf{0}\}$.
The coupling step is a bijection on $(\Ztwo)^{16}$---no information is lost.}

\paragraph{Eigenvalue table.}
Table~\ref{tab:eigenvalues} lists all 16 eigenvalue magnitudes.  All modes
satisfy $|\lambda_k| > 1$: every Fourier mode is amplified, none contracted.

\begin{table}[h]
\centering
\caption{Eigenvalue magnitudes $|\lambda_k| = |p(\omega^k)|$ for distances $\{1,3,7,11\}$.}
\label{tab:eigenvalues}
\begin{tabular}{@{}cc|cc@{}}
\toprule
$k$ & $|\lambda_k|$ & $k$ & $|\lambda_k|$ \\
\midrule
0 & 5.000 & 8 & 3.000 \\
1 & \textbf{1.259} & 9 & \textbf{1.259} \\
2 & 1.732 & 10 & 1.732 \\
3 & 2.101 & 11 & 2.101 \\
4 & 2.236 & 12 & 2.236 \\
5 & 2.101 & 13 & 2.101 \\
6 & 1.732 & 14 & 1.732 \\
7 & \textbf{1.259} & 15 & \textbf{1.259} \\
\bottomrule
\end{tabular}
\end{table}

%% ============================================================
%% 5. SECURITY ANALYSIS
%% ============================================================
\section{Security Analysis}
\label{sec:security}

\subsection{Permutation Properties}

\proof{The CML-Sponge permutation is a bijection on $(\Ztwo)^{16} \times \Ztwo$.
Each of the four round steps is individually bijective:}
\begin{itemize}[nosep]
\item Step~1 (counter injection): translation by a constant---invertible.
\item Step~2 (Cat Map): $\det(\catmap) = 1$---invertible over any ring.
\item Step~3 (coupling): $\det(\coupmat) = -33075$ (odd)---invertible over~$\Ztwo$.
\item Step~4 (multiplicative mixing): $(s[2k] \mathbin{|} 1)$ is always odd---invertible
  multiplier in~$\Ztwo$.
\end{itemize}

\proof{The effective capacity is the full $c = 512$~bits.
The coupling matrix has trivial kernel over~$\Ztwo$, so no state-space dimension
is collapsed.  No capacity correction is needed.}

\proof{The sponge security bound, conditional on the permutation being ideal, is:
\[
  \mathrm{Adv} \leq \frac{q^2}{2^{512}}.
\]
For $q = 2^{128}$ queries (128-bit security target):
$\mathrm{Adv} \leq 2^{256}/2^{512} = 2^{-256}$, which is negligible.}

\proof{Full 16-site diffusion occurs in exactly 2 rounds with distances $\{1,3,7,11\}$.
Verified by symbolic propagation: after 1~round, each site is influenced by
5 source sites; after 2~rounds, all 16~sites are mutually dependent.
The 8-round permutation provides a $4\times$ diffusion margin.}

\subsection{The Core Security Conjecture}

The entire security argument rests on one conjecture:

\begin{quote}
\conjecture{The CML-Sponge permutation $\pi$ is computationally indistinguishable
from a uniformly random permutation on $\{0,1\}^{1024}$.}
\end{quote}

\paragraph{What is NOT proven.}
\begin{itemize}[nosep]
\item No differential cryptanalysis of the CML round has been performed.
  The DDT of Step~4 has not been computed.
\item No linear cryptanalysis.  The LAT of Step~4 has not been computed.
\item No algebraic degree analysis.  The growth of algebraic degree through
  8~rounds has not been formally bounded.
\item No formal reduction from the permutation conjecture to a standard
  hardness assumption.
\end{itemize}

\paragraph{What IS proven.}
\begin{itemize}[nosep]
\item Bijectivity of every round step (proof above).
\item Coupling matrix: full rank, trivial kernel, all eigenvalues amplifying
  ($|\lambda_k| \geq 1.259$).
\item Cat Map: Anosov hyperbolicity with Lyapunov exponent $\approx 0.9624$.
\item No complement symmetry: verified algebraically and by test vectors.
\item Full 16-site diffusion in 2~rounds.
\end{itemize}

\paragraph{What is supported empirically.}
\begin{itemize}[nosep]
\item PractRand at 1~TB: 397~tests, zero persistent anomalies (Section~\ref{sec:validation}).
\item Raw permutation (no Mix13) at 32~GB: 325~tests, zero failures.
\item Reduced-round analysis: $4\times$ margin over statistical distinguisher threshold.
\end{itemize}

\subsection{Cryptanalysis Record}

Table~\ref{tab:attacks} summarizes the outcomes of all considered attacks.

\begin{table}[h]
\centering
\caption{Cryptanalysis attempts.  ``Closed'' means the attack was analyzed and
found to be non-viable.  ``Open'' means the analysis is incomplete.}
\label{tab:attacks}
\small
\begin{tabular}{@{}clll@{}}
\toprule
\# & Attack & Outcome & Status \\
\midrule
1 & Generic preimage (sponge) & No attack found & Closed (conditional) \\
2 & Coupling null vector & \textbf{Found \& fixed} ($\{1,7,8\}$) & \textbf{Closed (v10)} \\
3 & Statistical distinguisher & No anomaly at 1~TB & Closed (empirical) \\
4 & Related-key/IV & No attack found & Open \\
5 & Differential (1-round) & Not formally analyzed & Open \\
6 & Eigenvalue basin & All $|\lambda_k| > 1$ & Closed (v10) \\
7 & Diffusion horizon (1-round) & Confirmed at 1~round; closed at 2+ & Closed \\
8 & Even det (non-invertible) & \textbf{Found \& fixed} ($\{1,5,11\}$) & \textbf{Closed (v10)} \\
\bottomrule
\end{tabular}
\end{table}

\paragraph{Attack~2: Coupling singularity (the methodological contribution).}
The discovery that $p(-1) = 0$ for distances $\{1,7,8\}$ was made through
geometric analysis of the CML lattice dynamics.  The null vector
$\mathbf{v} = (1, -1, 1, -1, \ldots)$ makes the coupling step singular:
two states differing by~$\mathbf{v}$ produce identical output.
This \emph{invalidates the sponge security proof} by reducing the effective
state dimension.  The fix---a systematic search over all 1365 distance
combinations (Section~\ref{sec:coupling})---demonstrates that principled
cryptographic design can discover and resolve structural flaws that would
be invisible to statistical testing alone.  PractRand cannot detect this
singularity: the null vector is a single algebraic direction in a
$2^{1024}$-element state space.

\paragraph{Attack~8: Even determinant.}
The intermediate design $\{1,5,11\}$ fixed the $p(-1) = 0$ singularity
but introduced a subtler problem: $\det(\coupmat) = -1088$ is even,
giving a 4-element kernel and 2-bit capacity loss.  A 3-term coupling was
evaluated and rejected (Section~\ref{sec:coupling}).  The 5-term architecture
resolves both issues simultaneously: odd determinant, trivial kernel,
full capacity.

\subsection{The Nonlinearity Question}
\label{sec:nonlinearity}

We address the most important open question directly.

Arnold's Cat Map is a \emph{linear} map over~$\Ztwo$.  It is the matrix
$\begin{psmallmatrix}1&1\\1&2\end{psmallmatrix}$ acting by wrapping addition---degree~1 in both variables.  The coupling step (Step~3) is also linear (a
circulant matrix over~$\Ztwo$).  Steps~1 and~2 together form an \emph{affine}
map.

\textbf{Step~4 is the sole source of nonlinearity.}
$s[2k+1] \gets s[2k+1] \boxtimes (s[2k] \mathbin{|} 1)$ is degree-2 over~$\Ztwo$.
It acts on 8 of the 16~sites (the odd-indexed ones).  Over 8~rounds, the
algebraic degree grows roughly as $2^r$ for the affected words, but the
even-indexed words are only indirectly affected through the coupling step.

\paragraph{Comparison.}
In AES, the S-box (GF$(2^8)$ inversion) provides maximum algebraic degree~254
per application, applied to \emph{all} 16~state bytes.  In Keccak, the $\chi$
step provides degree-2 nonlinearity applied to all 1600~bits.  In ChaCha20,
the nonlinear operation (addition mod~$2^{32}$) is interleaved densely in every
quarter-round.  CATWALK's nonlinear density is lower: one multiplicative step
per round on half the state.

\paragraph{Why this is an open problem, not a known weakness.}
\begin{enumerate}[nosep]
\item No algebraic attack on the CML-Sponge permutation has been demonstrated.
  The linearity of Steps~1--3 does not automatically imply vulnerability;
  it means that the security depends critically on Step~4's interaction with
  the linear layers, which has not been analyzed.
\item The reduced-round analysis (Section~\ref{sec:validation}) provides
  empirical evidence: the \emph{raw} permutation (no Mix13) achieves
  statistical indistinguishability at 2~rounds.  If the nonlinearity were
  insufficient, we would expect to see statistical weaknesses persisting
  to higher round counts.
\item The Cat Map's linearity is over~$\Ztwo$, not over~$\mathrm{GF}(2)$.
  The wrapping-addition arithmetic means that carry propagation introduces
  effective nonlinearity at the bit level---a $k$-bit wrapping addition has
  algebraic degree~$k$ over~$\mathrm{GF}(2)$.  The cryptographic significance
  of this bit-level nonlinearity is not well understood.
\end{enumerate}

\paragraph{What future analysis should determine.}
\begin{itemize}[nosep]
\item Compute the differential distribution table of Step~4:
  $\Pr[\Delta y' \mid \Delta x, \Delta y]$ where $y' = y \boxtimes (x \mathbin{|} 1)$.
\item Compute the linear approximation table of Step~4.
\item Bound the algebraic degree of the full permutation output as a function
  of round count.
\item Determine whether the affine structure of Steps~1--3 permits efficient
  linearization attacks (Gr\"obner bases, MQ solvers).
\end{itemize}

\subsection{Geometric Foundations}

\paragraph{The state space.}
The CML-Sponge state can be viewed as a point in $(\mathbb{T}^2_{64})^8$,
where $\mathbb{T}^2_{64} = (\Ztwo)^2$ is the discrete torus.  Each Cat Map
pair lives on one copy of~$\mathbb{T}^2_{64}$, and the coupling step
provides inter-torus communication.

\paragraph{Anosov dynamics.}
On the real torus~$\mathbb{T}^2$, the Cat Map is an Anosov diffeomorphism:
it has a global hyperbolic structure with stable and unstable manifolds
filling the torus densely.  The Lyapunov exponent $\Lambda \approx 0.9624$
quantifies the exponential divergence rate.  On the discrete torus $\mathbb{T}^2_{64}$,
the map retains the algebraic structure (bijectivity, eigenvalue ratio) but
has finite period.  The cryptographic significance is that any two
``nearby'' states (in the Hamming metric) diverge rapidly through the
Cat Map---a form of the avalanche effect grounded in proven dynamical
systems theory rather than empirical observation.

\paragraph{Symplectic structure.}
$\catmap \in \mathrm{Sp}(2, \mathbb{Z})$ since $\catmap^T J \catmap = J$
where $J = \begin{psmallmatrix}0&1\\-1&0\end{psmallmatrix}$.
The symplectic property ensures area preservation (no contraction or
expansion of phase-space volume), which prevents state collapse.

%% ============================================================
%% 6. STATISTICAL VALIDATION
%% ============================================================
\section{Statistical Validation}
\label{sec:validation}

\subsection{Methodology}

All statistical testing uses PractRand~0.95~\cite{practrand} in \texttt{stdin64}
mode with the full core test suite (64-bit folding).  PractRand applies a
battery of up to 397~distinct statistical tests at the 1~TB checkpoint,
testing for correlations, bit-level biases, gap distributions, and other
non-random patterns.  Its four-level severity scale---``unusual'' ($p < 10^{-3}$),
``suspicious'' ($p < 10^{-5}$), ``very suspicious'' ($p < 10^{-7}$), and
``FAIL''---allows precise characterization of anomalies.

\subsection{Main Validation: 1~TB}

The v10 construction ($\{1,3,7,11\}$ coupling) was tested to 1~TB per seed.
Key and IV for each seed are derived via BLAKE3 from independent context strings.

\begin{table}[h]
\centering
\caption{PractRand v10 validation at 1~TB (seeds 0 and 1 complete).}
\label{tab:practrand-1tb}
\begin{tabular}{@{}clccl@{}}
\toprule
Seed & Key type & Tests at 1~TB & Anomalies (final) & Transient \\
\midrule
0 & BLAKE3-derived & 397 & 1 unusual$^*$ & 0 \\
1 & BLAKE3-derived & 397 & 0 & 1 unusual at 256~MB \\
2--255 & various & \multicolumn{3}{c}{In progress} \\
\bottomrule
\end{tabular}
\\[4pt]
{\footnotesize $^*$\texttt{[Low1/64]mod3n(5):(1,9-0)} $R = +8.3$, $p = 1.0 \times 10^{-4}$---
appeared only at the 1~TB checkpoint, consistent with the expected
Poisson rate of $\sim$0.4 unusual events per checkpoint at 397~tests.}
\end{table}

Prior to the v10 coupling design, an intermediate design ($\{1,5,11\}$, 4-term
coupling) was validated at 256~GB on 5~seeds (0, 1, 2, 254, 255), each passing
369~tests with zero persistent anomalies.  Those results are superseded by the
v10 validation reported here.

\subsection{Raw Permutation Validation}

To determine whether the Stafford Mix13 output finalizer is defense-in-depth
or load-bearing, the raw permutation output (before Mix13) was tested with
PractRand to 32~GB on three seeds.

\begin{table}[h]
\centering
\caption{Raw permutation PractRand results (no Mix13 finalizer).}
\label{tab:practrand-raw}
\begin{tabular}{@{}ccccc@{}}
\toprule
Seed & Key type & Tests at 32~GB & Failures & Transient unusual \\
\midrule
0 & BLAKE3-derived & 325 & 0 & 1 (at 1~GB, gone by 2~GB) \\
254 & $\texttt{0x00} \times 32$ & 325 & 0 & 1 (at 8--16~GB, gone by 32~GB) \\
255 & $\texttt{0xFF} \times 32$ & 325 & 0 & 1 (at 2~GB, gone by 4~GB) \\
\midrule
\textbf{Total} & & \textbf{975} & \textbf{0} & 3 (all transient) \\
\bottomrule
\end{tabular}
\end{table}

\empirical{Mix13 is defense-in-depth, not load-bearing.  The CML-Sponge
permutation produces statistically strong output at the raw lattice level.}

\subsection{Reduced-Round Analysis}

Rounds~1--8 were tested in both with-Mix13 and raw (no-Mix13) configurations,
each to 4~GB (Table~\ref{tab:reduced-round}).

\begin{table}[h]
\centering
\caption{Reduced-round PractRand results.  ``Pass'' = 277 tests at 4~GB, zero failures.}
\label{tab:reduced-round}
\begin{tabular}{@{}ccc@{}}
\toprule
Rounds & With Mix13 & Raw (no Mix13) \\
\midrule
1 & Pass & \textbf{FAIL} at 512~MB \\
2 & Pass & Pass \\
3 & Pass & Pass \\
4 & Pass & Pass \\
5 & Pass & Pass \\
6 & Pass & Pass (3 transient unusual) \\
7 & Pass & Pass \\
8 & Pass & Pass \\
\bottomrule
\end{tabular}
\end{table}

The 1-round raw failure is catastrophic: \texttt{[Low1/64]Gap-16:A}
$R = +189.2$, $p = 4.6 \times 10^{-162}$---extreme bias in the
least-significant-bit gap distribution, as expected from a single CML round
that provides only local pair-wise mixing without full 16-site diffusion.
At 2~rounds, full diffusion eliminates the distinguisher completely.

\paragraph{Security margin.}
\empirical{The raw permutation achieves statistical indistinguishability at
2~rounds.  The 8-round design provides a $4\times$ conservative margin
($8/2 = 4$).}  With Mix13, even 1~round passes (effective margin $\geq 8\times$),
but the conservative figure uses the raw permutation.

\begin{table}[h]
\centering
\caption{Security margin comparison with established ciphers.}
\label{tab:margin-comparison}
\begin{tabular}{@{}lccc@{}}
\toprule
Cipher & Design rounds & Distinguisher at & Margin \\
\midrule
ChaCha20~\cite{bernstein2008chacha} & 20 & 7 rounds & $2.9\times$ \\
AES-128~\cite{daemen2001aes} & 10 & 6 rounds & $1.7\times$ \\
CATWALK v10 & 8 & 2 rounds (raw) & $4.0\times$ \\
\bottomrule
\end{tabular}
\\[4pt]
{\footnotesize Note: ChaCha20 and AES margins are based on best published
reduced-round distinguishers.  CATWALK's margin is measured against statistical
distinguishers only; algebraic distinguishers have not been analyzed.}
\end{table}

\paragraph{Honest caveat.}
\empirical{PractRand validates output stream quality, not cryptographic security.
A generator that passes PractRand may still be broken by algebraic or structural
attacks.  The Mix13 output finalizer may mask algebraic structure that would be
exploitable by a cryptanalyst but invisible to statistical tests.  The
reduced-round margin should therefore be interpreted as a lower bound on the
true security margin, not a definitive characterization.}

%% ============================================================
%% 7. PERFORMANCE
%% ============================================================
\section{Performance}
\label{sec:performance}

All benchmarks were measured with Criterion~0.5 on Windows~10 IoT Enterprise
(x86-64, \todo{insert CPU model}),%
\space Rust stable with release profile (LTO, single codegen unit, strip).

\begin{table}[h]
\centering
\caption{CATWALK performance benchmarks.}
\label{tab:performance}
\begin{tabular}{@{}lrr@{}}
\toprule
Operation & Time & Throughput \\
\midrule
Permutation (8 rounds) & 137.6~ns & --- \\
Keystream (64 bytes) & 169.8~ns & 360~MiB/s \\
AEAD encrypt (1~KB) & 7.21~$\mu$s & 135~MiB/s \\
AEAD encrypt (64~KB) & 311.0~$\mu$s & 201~MiB/s \\
AEAD encrypt (1~MB) & 6.05~ms & 165~MiB/s \\
AEAD decrypt (1~KB) & 8.71~$\mu$s & 112~MiB/s \\
AEAD decrypt (64~KB) & 317.3~$\mu$s & 197~MiB/s \\
KDF (min params: 64~MB, $t\!=\!2$) & 93.4~ms & --- \\
KDF (production: 256~MB, $t\!=\!4$) & $\sim$1.5~s & --- \\
\bottomrule
\end{tabular}
\end{table}

\paragraph{Permutation cost.}
Each permutation applies 8~CML rounds over the 1024-bit state in 137.6~ns,
yielding a raw permutation ceiling of $64\text{ bytes} / 137.6\text{ ns} \approx 445$~MiB/s.

\paragraph{Streaming throughput.}
The AEAD plateau is $\sim$200~MiB/s at 64~KB payloads.  Each 64-byte ciphertext
block requires two permutations (one for squeeze, one for absorb), giving a
theoretical ceiling of $64 / (2 \times 137.6) \approx 222$~MiB/s, consistent
with the measured 200~MiB/s.

\paragraph{Comparison.}
ChaCha20-Poly1305 achieves 1--3~GiB/s on the same class of hardware (5--15$\times$
faster), leveraging SIMD-friendly quarter-round operations.  CATWALK's higher
per-byte cost is inherent to the CML architecture: the 16-site coupling step
involves 80~wrapping additions per round (5~per site $\times$ 16~sites), and the
AEAD absorb step doubles the permutation count per output block.

\paragraph{KDF dominance.}
For files smaller than a few hundred megabytes, the Argon2id KDF ($\sim$1.5~s
at production parameters) completely dominates the encryption time.  This is
intentional: the KDF cost makes offline password-guessing attacks expensive.

\paragraph{SIMD optimization potential.}
The multi-pass coupling implementation (5~sequential loops, one per distance)
was specifically designed for auto-vectorization.  The compiler generates
SIMD loads for each rotated-offset pass.  Further optimization with explicit
AVX2 intrinsics is future work.

%% ============================================================
%% 8. RELATED WORK
%% ============================================================
\section{Related Work}
\label{sec:related}

\paragraph{Chaos-based cryptography.}
Baptista~\cite{baptista1998cryptography} proposed the first chaos-based
encryption scheme, using the logistic map's sensitivity to initial conditions
as a keying mechanism.  Kocarev~\cite{kocarev2001chaos} surveyed the field
and identified key challenges: discretization artifacts, short periods on
finite-precision representations, and the gap between continuous-time chaos
theory and discrete cryptographic requirements.
Li~et~al.~\cite{li2001problems} provided the most influential critique,
demonstrating concrete attacks on multiple chaos ciphers based on
floating-point precision loss, complement symmetry, and weak mixing.
CATWALK addresses each of Li~et~al.'s concerns: exact integer arithmetic
eliminates precision loss, proven Anosov structure replaces assumed chaos,
and the sponge framework provides a formal security model.

\paragraph{Coupled Map Lattices in cryptography.}
Kaneko~\cite{kaneko1984cml,kaneko1985cml} introduced CMLs as models of
spatiotemporal chaos.  Their application to cryptography has been explored
by several authors, typically using logistic or tent maps as the local function
with diffusive coupling.  These proposals generally suffer from the same
floating-point issues as other chaos ciphers, and none embed the CML in
a sponge framework with formal capacity analysis.  CATWALK's use of
Arnold's Cat Map (integer-exact, provably hyperbolic) as the local map,
combined with the sponge wrapper, distinguishes it from prior CML proposals.

\paragraph{Sponge constructions.}
The sponge construction was introduced by Bertoni~et~al.~\cite{bertoni2007sponge}
and its indifferentiability from a random oracle was proved
in~\cite{bertoni2008indifferentiability}.  Keccak~\cite{keccak2011}, selected
as SHA-3, is the canonical instantiation.  CATWALK follows the same framework
but replaces the Keccak-$f$ permutation with the CML-Sponge permutation.
The sponge security bound $q^2/2^c$ applies identically, conditional on the
permutation's quality.

\paragraph{Arnold's Cat Map in image encryption.}
Fridrich~\cite{fridrich1998symmetric} used the Cat Map for pixel-position
scrambling in image encryption, exploiting the map's mixing property on the
torus.  Chen~et~al.~\cite{chen2004symmetric} extended this to 3D Cat Maps.
These applications use the Cat Map as a position permutation (shuffling pixel
locations), not as a state-evolution primitive within a sponge.  CATWALK uses
the Cat Map in a fundamentally different way: as the local map in a CML that
serves as the sponge permutation's core mixing step.

\paragraph{ARX stream ciphers.}
Salsa20~\cite{bernstein2005salsa20} and ChaCha20~\cite{bernstein2008chacha}
are the dominant modern stream ciphers, built on the ARX (add-rotate-XOR)
paradigm.  Their quarter-round structure provides dense nonlinear mixing
through the interaction of addition mod~$2^{32}$ and XOR.  CATWALK's round
function uses a similar vocabulary of operations (wrapping addition,
multiplication, rotation) but in a different architecture: the CML topology
provides the diffusion mechanism rather than a fixed-wiring permutation.
ChaCha20 achieves 5--15$\times$ higher throughput and has received extensive
cryptanalysis; CATWALK's advantage is the richer mathematical structure
of the CML dynamics.

\paragraph{AES and block ciphers.}
AES~\cite{daemen2001aes,fips197} uses the SPN (substitution-permutation network)
paradigm with a high-degree S-box (GF$(2^8)$ inversion, algebraic degree~254).
This provides maximum nonlinear density per round.  CATWALK's lower nonlinear
density (Section~\ref{sec:nonlinearity}) is compensated by the larger state
(1024 vs.\ 128~bits) and higher round count relative to the diffusion minimum
($4\times$ vs.\ AES's $1.7\times$).  Whether this trade-off is sufficient for
cryptographic security is the central open question.

\paragraph{Password-based key derivation.}
Argon2~\cite{rfc9106}, the winner of the Password Hashing Competition, provides
memory-hard key derivation.  CATWALK uses Argon2id (hybrid data-dependent and
data-independent passes) at 256~MB / 4~iterations, followed by BLAKE3~\cite{blake3}
domain derivation.  The decryption-side parameter floor prevents KDF downgrade
attacks via crafted file headers.

%% ============================================================
%% 9. OPEN PROBLEMS
%% ============================================================
\section{Open Problems}
\label{sec:openproblems}

We enumerate the unresolved questions in order of estimated importance.
Each represents an opportunity for independent cryptanalysis.

\begin{enumerate}

\item \textbf{Formal pseudorandomness reduction.}
  \emph{Problem:} Prove (or disprove) that the CML-Sponge permutation is
  a pseudorandom permutation under a standard hardness assumption.
  \emph{Why it matters:} Without this, the security argument is conditional
  on an unverified conjecture.
  \emph{What closing it requires:} A complexity-theoretic reduction, or a
  concrete attack demonstrating the conjecture is false.
  \emph{Status:} Completely open.

\item \textbf{Differential and linear cryptanalysis of the round function.}
  \emph{Problem:} Compute or bound the maximum differential probability and
  linear correlation of the CML round.
  \emph{Why it matters:} Steps~1--3 are affine; differential/linear propagation
  through them is deterministic.  The security depends entirely on Step~4's
  resistance.
  \emph{What closing it requires:} Computation of the DDT and/or LAT for
  Step~4, then bounding multi-round characteristics using the coupling
  structure for diffusion.
  \emph{Status:} Not attempted.

\item \textbf{Algebraic degree analysis.}
  \emph{Problem:} Compute the algebraic degree of the permutation output as a
  function of round count, particularly for the even-indexed sites that are
  only indirectly affected by Step~4.
  \emph{Why it matters:} Low algebraic degree after 8~rounds would indicate
  vulnerability to algebraic attacks (Gr\"obner bases, MQ solvers).
  \emph{What closing it requires:} Symbolic computation tracking the degree
  growth through the composition of affine and multiplicative steps.
  \emph{Status:} Not attempted.

\item \textbf{Period of the Cat Map on $(\Ztwo)^2$.}
  \emph{Problem:} Compute or tightly bound the period of
  $\begin{psmallmatrix}1&1\\1&2\end{psmallmatrix}$ acting on $(\Ztwo)^2$.
  \emph{Why it matters:} A short period would limit the keystream length.
  In practice, the Weyl counter injection ensures each round operates
  at a different offset, making the bare period largely irrelevant.
  \emph{What closing it requires:} Number-theoretic analysis of the matrix
  order modulo~$2^{64}$.
  \emph{Status:} Not computed.  \conjecture{The period is astronomically large.}

\item \textbf{Independent cryptanalysis.}
  \emph{Problem:} Subject the CML-Sponge permutation to analysis by
  researchers with expertise in symmetric cryptanalysis.
  \emph{Why it matters:} Self-review is inherently limited.
  \emph{What closing it requires:} External engagement.
  \emph{Status:} This paper is intended to facilitate exactly this.

\item \textbf{Mixing proof: exponential correlation decay.}
  \emph{Problem:} Prove that spatial correlations in the CML lattice decay
  exponentially with round count.
  \emph{Why it matters:} Would provide a theoretical foundation for the
  empirical observation that 2~rounds suffice for statistical indistinguishability.
  \emph{What closing it requires:} Transfer operator analysis of the CML
  dynamics over~$\Ztwo$, extending results from continuous CML theory.
  \emph{Status:} Not attempted.

\item \textbf{5-term coupling optimality.}
  \emph{Problem:} Determine whether the 5-term architecture with
  $\{1,3,7,11\}$ is optimal among all feasible coupling designs for this
  lattice size, or whether alternative structures (e.g., 6-term, non-uniform
  coefficients, non-circulant topologies) offer stronger security guarantees.
  \emph{Why it matters:} The exhaustive search covered only unit-coefficient
  circulant designs.
  \emph{What closing it requires:} Extended search or proof of optimality.
  \emph{Status:} The 5-term architecture is the best among unit-coefficient
  4-distance circulants; other design spaces are unexplored.

\end{enumerate}

%% ============================================================
%% 10. CONCLUSION
%% ============================================================
\section{Conclusion}
\label{sec:conclusion}

We have presented CATWALK, a research-grade authenticated encryption scheme
whose core permutation is a Coupled Map Lattice wrapped in a standard sponge
framework.  The construction is specified precisely enough for independent
reimplementation (Section~\ref{sec:construction}), and every security claim
is labeled as \textsc{[Proof]}, \textsc{[Conjecture]}, or \textsc{[Empirical]}.

The central methodological contribution is the design-by-analysis process
that discovered and resolved a structural singularity in the coupling topology
(Section~\ref{sec:coupling}).  Geometric analysis revealed the flaw;
exhaustive polynomial root analysis fixed it.  This process---where the
analysis tool and the design tool are the same mathematical framework---is,
we believe, a model for principled cryptographic design.

CATWALK is a research cipher.  It has not received independent cryptanalysis.
The core security conjecture---that the CML-Sponge permutation is
indistinguishable from random---remains unproven.  The nonlinearity question
(Section~\ref{sec:nonlinearity}) is the highest-priority open problem.
Statistical validation (1~TB PractRand, raw permutation testing, $4\times$
reduced-round margin) provides empirical support but is not a substitute for
formal analysis.

We explicitly invite the cryptographic community to scrutinize this
construction.  The open problems enumerated in Section~\ref{sec:openproblems}
represent concrete opportunities for analysis.  A successful attack would be
as valuable as a successful defense: it would clarify the boundary between
chaos-theoretic intuition and cryptographic reality.

The complete reference implementation (Rust, $\sim$530 lines for the core
primitive) and all documentation are publicly available.

\bibliographystyle{splncs04}
\bibliography{refs}

\end{document}
